\section{Introduction}
\label{sec:intro}

Parton distribution functions (PDFs) describe the inner dynamics of partons inside a hadron \cite{Feynman:102074}.
They have a non-pertubative nature and, thus, they can not be computed in perturbation theory. Lattice QCD is an
ideal formulation to study the PDFs from first principles, in large scale simulations.
However, PDFs are usually defined on the light cone, which poses a problem for the standard Euclidean formulation.
Hence,  hadron structure calculations in lattice QCD are related to other quantities  that are accessible in a Euclidean
spacetime. This led to a long history of investigations of Mellin moments of PDFs and nucleon form factors (see
\cite{Constantinou:2014tga, Constantinou:2015agp,Alexandrou:2015yqa, Alexandrou:2015xts, Syritsyn:2014saa}
for recent reviews). In practice, there are severe limitations in the reconstruction of the PDFs mainly due to the small signal-to-noise ratio for the
high moments. In addition,  there are inevitable problems with power divergent mixings with lower dimensional operators.
Therefore, the task of reconstructing the PDFs from their moments is practically unfeasible.

\medskip
{\textit{Ab initio}} evaluations of PDFs are of high importance as they would be a stringent test of non-perturbative
aspects of QCD. The fact that a calculation of the PDFs from first principles is missing is a pressing problem that
prevents a deeper understanding of the nucleon structure.
Our current knowledge on the PDFs relies on phenomenological fits from experimental data using perturbation theory.
These parameterized PDFs serve, for example, as input for the computation of cross sections used for
presently running colliders, most notably the LHC, and also to plan future collider experiments which are themselves
tests of the Standard Model. The parameterizations are, however, not without ambiguities~\cite{Jimenez-Delgado:2013sma}.
In addition, there are kinematical regions that are not experimentally accessible.
The large Bjorken $x$ region is one of them, with large uncertainties most dramatically seen in the down quark
distributions \cite{Accardi:2016qay,Alekhin:2017kpj}. The transversity PDF is yet another example of a PDF that is
only poorly constrained by phenomenology.

\medskip
A pioneering method for a direct computation has been suggested by X.~Ji \cite{Ji:2013dva}.
In this approach, instead of matrix elements defined on the light cone, one calculates matrix elements
of fermion operators including a finite-length Wilson line, and whose Fourier transform defines the so-called
quasi-PDFs. This is achieved by taking the Wilson line in a purely spatial direction, conventionally chosen
to be the $z$-direction, instead of the $+$-direction on the light cone.
In terms of hadron kinematics, the nucleon momentum is usually taken along this spatial direction.
At large, but finite momenta, the quasi-PDFs can then be related to the light-front
PDFs via a matching procedure~\cite{Xiong:2013bka,Chen:2016fxx}. Ji's approach has already been tested in
Refs.~\cite{Lin:2014zya,Gamberg:2014zwa,Alexandrou:2015rja,Chen:2016utp,Alexandrou:2016jqi} and all these results
are promising, i.e.\ they give a correct shape of the PDFs after the matching procedure. Certain properties of quasi-PDFs,
like the nucleon mass dependence and target mass effects, have also been analyzed via their relation with transverse momentum
dependent distribution functions (TMDs)~\cite{Radyushkin:2017ffo,Radyushkin:2016hsy}.
Note also the appearance of a related approach, pseudo-PDFs, which is a different generalization of light-cone PDFs to finite nucleon momenta \cite{Radyushkin:2017cyf,Orginos:2017kos}.
Refs.~\cite{Carlson:2017gpk,Briceno:2017cpo} discussed the role of the Euclidean signature in the computation of quasi-PDFs, as compared to light-front PDFs, and the latter reference proved that matrix elements obtained from Euclidean lattice QCD are identical to those obtained using the LSZ reduction formula in Minkowski space.
Nevertheless, the matrix elements of quasi-PDFs contain divergences that need to be eliminated via renormalization in order to obtain
meaningful results that can be compared to the physical PDFs.

\medskip
To date, all works on the quasi-PDFs only considered the bare matrix elements, as the renormalization process is
highly non-trivial and was not addressed until recently. In particular, new complications arise in the
renormalization of the Wilson line operators, compared to the local operators. For one, in addition to the
logarithmic divergences there is a linear divergence~\cite{Dotsenko:1979wR} with respect to the lattice regulator,
$a$, that prohibits one to take the continuum limit prior to its elimination. To one-loop level in perturbation theory the
divergence is manifesting itself as a linear divergence, computed in Ref.~\cite{Constantinou:2017sej} for a variety
of fermion and gluon actions. However, it is of utmost importance to extract the power divergence non-perturbatively,
which is one of the goals of this work.  Another feature of these operators that brings in new complications is the fact
that certain choices of the Dirac structure exhibit mixing~\cite{Constantinou:2017sej}.

\medskip
To show that these matrix elements can be renormalized, in particular that the linear divergence associated with
the Wilson line can be eliminated, is of paramount importance. Without renormalization, the whole quasi-PDF strategy
is incomplete and unable to provide any useful information to the theoretical and experimental community. Some suggestions
for the elimination of the linear divergence via the static potential were proposed in Refs.~\cite{Ma:2014jla,Ishikawa:2016znu}.
In Ref.~\cite{Chen:2016fxx} a one loop calculation of the linear divergence has been made, and that motivated the definition of an
improved quasi-PDF that is free of power divergences. One has, nevertheless, to show that such a procedure can be done non-perturbatively.
Another method to extract the coefficient of the linear divergence using the nucleon matrix elements of the quasi-PDFs was also presented
in Ref.~\cite{Constantinou:2017sej}. An alternative technique to suppress the linear divergence was discussed in Ref.~\cite{Monahan:2016bvm}
utilizing the gradient flow. Very recently, two papers discussed the employment of the auxiliary field formalism for the renormalization of
quasi-PDFs~\cite{Ji:2017oey,Green:2017xeu}, where the Wilson line is replaced by a Green's function of the introduced auxiliary field.
The renormalizability of quasi-PDFs to all orders in perturbation theory was addressed in Ref.~\cite{Ishikawa:2017faj}.

\medskip
In this paper we propose, for the first time\footnote{After the submission of our work, the proposed renormalization programme was also applied in Ref.~\cite{Chen:2017mzz}.}, a concrete renormalization method of the quasi-PDFs in a fully non-perturbative manner.
We provide the prescription of the method and show examples of the renormalized matrix elements. We also discuss the
elimination of the mixing between the unpolarized quasi-PDF and the twist-3 scalar operator.
We employ the RI$'$ renormalization scheme~\cite{Martinelli:1994ty} and we convert the results to the
$\MSb$ scheme at a reference scale, $\bar\mu{=}2$ GeV, using the one-loop conversion factor computed in Ref.~\cite{Constantinou:2017sej}.
As a test case, we focus on the helicity quasi-PDF to demonstrate results, as it is free of mixing.

\medskip
The outline of the paper is as follows: In Section 2 we provide the theoretical setup related to the nucleon
matrix elements and the renormalization prescription in the presence and absence of mixing,
as well as the data on the conversion factor computed perturbatively.
Section 3 includes results on the
renormalization functions, a discussion on the systematic uncertainties in the $Z$-factors, renormalized matrix elements for the helicity case, as well as results of the matching to light-front PDFs. Finally we conclude and give our future directions.
